\section{Introduction}
\label{sec:intro}

\subsection{Motivation}

Machine learning systems are stochastic artefacts.
Even when the architecture, loss, and hyperparameters are fixed, training outcomes vary with the realisation of random draws used for initialisation, data ordering, regularisation masks, and augmentation~\cite{summers2021nondeterminism,zhuang2022randomness,pham2022quantifying}.
Optimisation methods themselves---stochastic gradient descent (SGD) and its adaptive variants, as well as population-based metaheuristics---are defined by how they inject and structure randomness into the search process~\cite{keskar2017largebatch,smith2018bayesian,zhu2019anisotropic}.

A separate, older literature asks a different question: when is a sequence of bits ``random enough''?
Statistical test batteries such as NIST~SP~800-22, Dieharder, and TestU01~\cite{bassham2010nist,lecuyer2007testu01} and, more recently, neural distinguishers~\cite{valle2024assessing} attempt to answer that question for cryptography and simulation.
The present review sits at the junction of these two traditions.
If an organisation holds \emph{very large binary random archives}---from high-quality PRNGs, hardware TRNGs, or QRNGs---can differences in measured bitstream quality be shown to matter for AI/ML task quality, training stability, or optimiser behaviour?

\subsection{Scope}

We review and critically assess work on:
\begin{enumerate}[leftmargin=*]
  \item definitions and measurement of RNG quality (classical batteries; entropy estimates; ML-based tests);
  \item substitution of PRNG / TRNG / QRNG sources into neural network training (initialisation, dropout, shuffling);
  \item the role of randomness and noise structure in first-order optimisers (SGD, Adam and relatives) and in metaheuristics (GA, PSO, DE);
  \item reproducibility, tooling nondeterminism, and reporting standards for seed-sensitive studies.
\end{enumerate}
We deliberately treat \emph{seed sensitivity} (different seeds of the same generator) and \emph{generator quality} (different statistical properties of the source) as distinct constructs; conflating them is a recurring methodological failure in the literature (\S\ref{sec:critique}).

\subsection{Contributions of this review}

\begin{enumerate}[leftmargin=*]
  \item A structured map of how randomness enters ML and optimisation pipelines, and which notions of ``quality'' are claimed to matter.
  \item A critical synthesis of conflicting empirical claims (especially QRNG vs.\ PRNG), with attention to statistical power, effect size, and confounds.
  \item An explicit gap analysis connecting cryptographic RNG evaluation, framework PRNG implementations, and optimiser noise theory.
  \item A research agenda for controlled experiments with large binary random bitstreams (\S\ref{sec:agenda}), aligned with the accompanying project plan.
\end{enumerate}

\subsection{Organisation}

Section~\ref{sec:background} fixes terminology and quality metrics.
Section~\ref{sec:rngml} reviews RNG substitution studies in ML.
Section~\ref{sec:opt} reviews randomness in stochastic and metaheuristic optimisation.
Section~\ref{sec:critique} offers a critical synthesis.
Section~\ref{sec:gaps} states research gaps.
Section~\ref{sec:agenda} outlines a concrete experimental programme.
